\vfill\pagebreak
\section{Boolean type and Boolean algebra}
\label{sec:types:boolean}

A Boolean holds one of two values, \token{true} or \token{false}. That is the
whole type, and it is enough to build every decision a program makes: which
branch to take, whether to loop again, whether to stop. Ymir calls it
\token{bool} and stores it in one byte, in which \token{0b00000000} is
\token{false} and every other pattern is \token{true}.

\subsection{Primary operators}

Boolean algebra is the arithmetic of truth values. Where numerical algebra has
four basic operations, this one has three -- AND, OR and NOT -- and every other
logical operation can be written out of those three alone.

AND, or conjunction, is $\land$ in algebra and \token{&&} in Ymir. It takes two
operands and yields \token{true} only when both are \token{true}. OR, or
disjunction, is $\lor$ and \token{||}: it yields \token{true} when at least one
operand is. NOT, or negation, is $\lnot$ and \token{!}, takes one operand, and
returns the other value. In the tables that follow, \bot{} stands for
\token{false} and \top{} for \token{true}.

\input{chapters/chapter2/figures/boolean_algebra_base_operators}
\vspace{-10pt}%
\subsection{Secondary operators}

Three further operators are common enough to have names of their own --
implication, bidirectional implication and exclusive disjunction -- but none of
them is new: each rewrites into the three basic ones.

\begin{itemize}
\vspace{-5pt}%
\item Implication: $x \rightarrow y = \lnot x \lor y$.
\vspace{-5pt}%
\item Bidirectional implication : $x \leftrightarrow y = (\lnot x \lor y) \land (x \lor \lnot y)$ 
\vspace{-5pt}
\item Exclusive disjunction (or XOR) : $x \oplus y = (x \land \lnot y) \lor (\lnot x \land y)$
\end{itemize}
\vspace{-5pt}%

\input{chapters/chapter2/figures/boolean_algebra_secondary}


The truth table above is what justifies those rewritings. Take implication. When
$x$ is true, the implication holds only if $y$ is true as well -- that is the
case everyday language recognises. When $x$ is false, the implication says
nothing about $y$ and holds either way, because implication runs one way only:
the truth of $y$ tells us nothing about $x$. Read as English, $\lnot x \lor y$
sounds like a different claim; row for row, it is the same one.

Ymir provides none of the three. The primary operators already express them, and
a language gains nothing by spelling them out.

\subsection{Laws on Boolean operators}

Boolean operations obey much the same laws as arithmetic ones, with $\lor$ in the
role of addition and $\land$ in the role of multiplication. Commutativity,
associativity and distributivity all carry over.

\textit{a)} Commutativity: the order of the operands does not affect the result of the
operation. This property applies to both AND and OR, for example $x \lor y =
y \lor x$.
\smallskip

\textit{b)} Associativity: the way operands are grouped does not affect the result of the
operation. This property applies to both AND and OR; for example, $x \land
(y \land z) = (x \land y) \land z$. However, associativity does not hold when
different operators are mixed. For instance, $x \land (y \lor z) \neq (x \land
y) \lor z$. To illustrate, let $x = \bot, y = \bot$ and $z = \top$, then
$x \land (y \lor z) = \bot$ while $(x \land y) \lor z = \top$.
\smallskip

\textit{c)} Distributivity: One operation can be spread across another. For instance
$x \land (y \lor z) = (x \land y) \lor (x \land z)$.

These laws are also where Boolean logic goes wrong most often. The analogy with
arithmetic extends to precedence -- \token{&&} binds tighter than \token{||},
just as multiplication binds tighter than addition -- and it is worth
parenthesising a Boolean expression even when the parentheses change nothing.

The classic mistake is to lean on commutativity and forget precedence, which
quietly regroups the operands. In the next listing \token{r1} and \token{r2} look
like rearrangements of one another and hold different values.

\begin{lstlisting}[style=coloredverbatim, label=lst:types:boolean_error, caption=Example of logic error when dealing with Boolean values, escapechar=@]
let x = false, y = false, z = true;

let r1 = x && y || z;
let r2 = (z || x) && y;

@\hb{assert(r1 == r2);}@ // Sorry, wrong. 
\end{lstlisting}

\subsection{Specificities of Boolean operators in Ymir}

Ymir compiles a Boolean operator into a branch, not into a calculation, and the
consequence is visible from the outside: the second operand is evaluated only if
it can still change the answer. A \token{false} on the left of \token{&&}, or a
\token{true} on the left of \token{||}, settles the result on its own, and the
right-hand side is skipped.

This buys two things. The skipped operand may be expensive -- a function call
rather than a variable -- so not evaluating it saves the work. More importantly,
it makes the left operand a guard for the right: later chapters lean on that
heavily.

In the listing below \token{foo} is never called. \token{x} is \token{true}, and
$\top \lor y$ is $\top$ whatever $y$ turns out to be. Swap the \token{||} on
line~11 for \token{&&} and \token{foo} does run -- unless \token{x} is also
changed to \token{false}, since $\bot \land y$ is $\bot$ regardless.

\begin{lstlisting}[style=coloredverbatim, caption=Example of boolean operation optimised out]
fn foo()-> bool {
  println("In foo");
  true
}

fn main()
  throws AssertError
{
  let mut x = true;

  assert(x || foo()); // not calling foo
  println("Exit");
}
\end{lstlisting}
\begin{lstlisting}[style=bashVerb, escapechar=@+]
@+\textcolor{teal!80}{alice@dev:\mtilde}@+$ gyc main.yr -o main
@+\textcolor{teal!80}{alice@dev:\mtilde}@+$ ./main
Exit
\end{lstlisting}


So the commutativity that Boolean algebra guarantees does not survive contact
with Ymir. The \emph{value} of \token{a && b} and \token{b && a} is the same; what
each of them evaluates on the way there is not.
